\documentclass[12pt,a4paper]{article}
\usepackage{amsmath,amssymb,amsthm}
\usepackage{hyperref}
\usepackage{geometry}
\geometry{margin=2.5cm}
\usepackage{enumitem}
\usepackage{booktabs}

\newtheorem{theorem}{Theorem}[section]
\newtheorem{lemma}[theorem]{Lemma}
\newtheorem{corollary}[theorem]{Corollary}
\newtheorem{proposition}[theorem]{Proposition}
\theoremstyle{definition}
\newtheorem{definition}[theorem]{Definition}
\theoremstyle{remark}
\newtheorem{remark}[theorem]{Remark}

\title{Fast Radio Bursts in Recognition Science:\\
Coherent Recognition Events in Magnetar Magnetospheres}

\author{Jonathan Washburn\\
Recognition Science Research Institute\\
Austin, Texas\\
\texttt{washburn.jonathan@gmail.com}}

\date{March 2026}

\begin{document}
\maketitle

\begin{abstract}
We analyze fast radio bursts (FRBs) within Recognition Science (RS).
FRBs---millisecond-duration radio transients with dispersion measures
indicating cosmological distances~\cite{Lorimer2007, Petroff2019}---are
interpreted as coherent recognition events in magnetar
magnetospheres.  The extreme brightness temperature
($T_b > 10^{35}\;\mathrm{K}$) requires a coherent emission
mechanism, which RS identifies as synchronized ledger updates
across ${\sim}\,10^{30}$ recognition events within a compact
emission region.  The characteristic frequency
$\nu \sim 1\;\mathrm{GHz}$ corresponds to the cyclotron frequency
in a $B \sim 10^{14}\;\mathrm{G}$ magnetar field, which sits on a
specific $\varphi$-rung of the magnetic field ladder.  The
dispersion measure--redshift relation
$\mathrm{DM} \propto \int n_e\,dl$ probes the intergalactic medium
and is consistent with RS cosmology.  Repeating FRBs are explained
by recurrent magnetospheric instabilities.
\end{abstract}

%% ============================================================
\section{Introduction}
\label{sec:intro}
%% ============================================================

Recognition Science~\cite{RS_Axioms} (RS) identifies FRBs as coherent
recognition events in magnetar magnetospheres, where the 8-tick epoch
synchronizes large numbers of ledger updates.

Since the discovery of the ``Lorimer burst''~\cite{Lorimer2007},
FRBs have become one of the most active areas of radio astronomy.
Key observational facts:
\begin{itemize}
  \item Duration: ${\sim}\,1$--$10\;\mathrm{ms}$.
  \item Dispersion measures: $\mathrm{DM} \sim 100$--$2000\;
  \mathrm{pc\,cm^{-3}}$, far exceeding Galactic contributions.
  \item Brightness temperature: $T_b > 10^{35}\;\mathrm{K}$,
  requiring coherent emission.
  \item Repetition: some sources repeat (e.g., FRB~121102,
  FRB~20201124A), others appear as one-off events.
  \item Association: FRB~200428 was detected from the Galactic
  magnetar SGR~1935+2154~\cite{CHIME2020}, confirming a magnetar
  origin for at least some FRBs.
\end{itemize}

%% ============================================================
\section{FRBs as Coherent Recognition Events}
\label{sec:mechanism}
%% ============================================================

\begin{definition}[Coherent Recognition Event]
A coherent recognition event is a synchronized update of
$N_{\mathrm{coh}}$ ledger entries within a single 8-tick epoch,
producing electromagnetic emission with power:
\begin{equation}
  P \propto N_{\mathrm{coh}}^2,
\end{equation}
rather than $P \propto N_{\mathrm{coh}}$ (incoherent emission).
\end{definition}

\begin{theorem}[FRB Coherence Number]
The observed FRB luminosity $L_{\mathrm{FRB}} \sim
10^{42}\;\mathrm{erg/s}$ requires:
\begin{equation}
  N_{\mathrm{coh}} \sim \sqrt{\frac{L_{\mathrm{FRB}}}{P_1}}
  \sim 10^{30},
\end{equation}
where $P_1 \sim 10^{-18}\;\mathrm{erg/s}$ is the single-particle
emission power.
\end{theorem}

\begin{proof}
For incoherent emission from $N$ particles, $P = N P_1$; for coherent
emission (all in phase), $P = N^2 P_1$ (constructive interference of
amplitudes).  From the observed FRB luminosity $L \sim 10^{42}$~erg/s
and a single cyclotron emission power $P_1 = (\sigma_T c) B^2/(8\pi)$
for one electron at $B \sim 10^2$ G (near emission radius), giving
$P_1 \sim 10^{-18}$ erg/s:
\[
  N_{\mathrm{coh}} = \sqrt{L/P_1} = \sqrt{10^{42}/10^{-18}} = 10^{30}.
\]
\end{proof}

\begin{remark}
This coherence number is enormous but achievable in a magnetar
magnetosphere where the magnetic field imposes a common axis and
frequency on all emitting particles.  The 8-tick epoch provides the
clock that synchronizes the recognition events.
\end{remark}

%% ============================================================
\section{Magnetar Connection}
\label{sec:magnetar}
%% ============================================================

\begin{theorem}[Cyclotron Frequency]
The characteristic FRB frequency $\nu \sim 1\;\mathrm{GHz}$
corresponds to the cyclotron frequency at a specific radius in the
magnetar magnetosphere:
\begin{equation}
  \nu_c = \frac{eB}{2\pi m_e c} \approx 2.8 \times
  10^6\;\mathrm{Hz} \times B(\mathrm{G}).
\end{equation}
For $\nu = 1\;\mathrm{GHz}$: $B \approx 360\;\mathrm{G}$, which
occurs at $r \approx R_{\mathrm{NS}} (B_s/360)^{1/3}$ in a dipole
field with surface field $B_s \sim 10^{14}\;\mathrm{G}$.
\end{theorem}

\begin{remark}
The emission radius for $\nu \sim 1\;\mathrm{GHz}$ is
$r \sim 10^4\;\mathrm{km} \sim 1000\,R_{\mathrm{NS}}$,
consistent with the magnetosphere.  The $\varphi$-ladder positions
the magnetar surface field $B_s \sim 10^{14}\;\mathrm{G}$ at a
specific rung, and the emission radius is determined by the dipole
geometry.
\end{remark}

%% ============================================================
\section{Dispersion Measure}
\label{sec:dm}
%% ============================================================

\begin{definition}[Dispersion Measure]
\begin{equation}
  \mathrm{DM} = \int_0^d n_e(l)\,dl,
\end{equation}
where $n_e$ is the free electron density along the line of sight.
\end{definition}

\begin{theorem}[DM--Redshift Relation]
For a source at redshift $z$:
\begin{equation}
  \mathrm{DM}_{\mathrm{IGM}} \approx 900\,z\;
  \mathrm{pc\,cm^{-3}},
\end{equation}
assuming the mean IGM electron density
$\bar{n}_e \approx 2 \times 10^{-7}\;\mathrm{cm^{-3}}$.
\end{theorem}

\begin{proof}
For a flat cosmology with Hubble radius $c/H_0 \approx 4285$~Mpc, the
comoving distance to redshift $z$ is $d \approx (c/H_0)z$ for $z \lesssim 1$.
The DM accumulates as $\mathrm{DM} = \bar{n}_e \cdot d$ (ignoring
$(1+z)$ evolution for simplicity):
\[
  \mathrm{DM} \approx (2 \times 10^{-7}\,\mathrm{cm}^{-3})
  \times (4285\,\mathrm{Mpc}) \times z
  \approx 900\,z\;\mathrm{pc\,cm}^{-3}.
\]
\end{proof}

\begin{remark}
This relation is consistent with RS cosmology (Hubble parameter
$H_0 \approx 71.4\;\mathrm{km/s/Mpc}$, dark energy fraction
$\Omega_\Lambda = 11/16 - \alpha/\pi \approx 0.685$).  FRBs serve
as cosmological probes of the IGM baryon content.
\end{remark}

%% ============================================================
\section{Repeating vs.\ One-Off FRBs}
\label{sec:repeaters}
%% ============================================================

\begin{theorem}[Repeater Mechanism]
Repeating FRBs arise from recurrent magnetospheric instabilities
in young magnetars ($\tau < 10^4\;\mathrm{yr}$).  Each burst
corresponds to a reconnection event that synchronizes a coherent
recognition event in a different part of the magnetosphere.
\end{theorem}

\begin{proof}[Structural argument]
Young magnetars have stronger spin-down torques and more rapidly
evolving magnetic field structures, creating recurrent twisted field
lines and reconnection events.  Each reconnection releases energy
sufficient to trigger a coherent recognition event (the reconnection
energy $\sim 10^{42}$--$10^{43}$ erg matches the FRB luminosity
integrated over duration).  In the RS framework, the 8-tick epoch
provides the synchronization clock: every 8-tick cycle in the
reconnection region that aligns $N_{\mathrm{coh}} \sim 10^{30}$
particles in phase produces one burst.  Recurrence occurs as the
field lines re-twist on the Alfvén crossing time $\tau_A \sim R/v_A$,
which is ${\sim}$ days to months for typical magnetar parameters.
\end{proof}

\begin{theorem}[One-Off Events]
Apparent one-off FRBs are either:
\begin{enumerate}
  \item Repeaters with recurrence times longer than the observation
  window.
  \item FRBs from older, less active magnetars with weaker and less
  frequent instabilities.
\end{enumerate}
RS predicts no fundamental distinction between repeaters and
non-repeaters---only differences in activity level.
\end{theorem}

\begin{proof}[Structural argument]
In RS, every FRB is a coherent recognition event triggered by a
magnetospheric instability.  The rate of such events is governed by
the magnetar's spin-down luminosity $\dot{E} \propto B^2/P^4$, where
$B$ is the surface field and $P$ is the spin period.  Old magnetars
have smaller $B$ and longer $P$, hence smaller $\dot{E}$, making
instabilities rarer.  An ``apparent'' one-off FRB is one whose
recurrence rate is so low ($\lesssim 1$ per year) that the probability
of repeating within a typical monitoring campaign is small.  The RS
prediction is that with sufficiently long monitoring, all FRBs will
repeat.  Non-repetition within any finite observation window cannot
prove a source is a true one-off.
\end{proof}

%% ============================================================
\section{Predictions}
\label{sec:predictions}
%% ============================================================

\begin{enumerate}
  \item \textbf{All FRBs are from magnetars}: No FRB should be
  associated with a non-magnetar source (e.g., black hole, white
  dwarf).

  \item \textbf{DM--$z$ relation}: FRBs with localized host galaxies
  should follow the $\mathrm{DM} \approx 900z$ relation, providing
  an independent measurement of the baryon density $\Omega_b$.

  \item \textbf{Polarization structure}: FRBs should show high
  linear polarization ($> 80\%$), consistent with the coherent
  emission mechanism and ordered magnetic field geometry.
\end{enumerate}

%% ============================================================
\section{Conclusion}
\label{sec:conclusion}
%% ============================================================

Fast radio bursts are coherent recognition events in magnetar
magnetospheres, requiring ${\sim}\,10^{30}$ synchronized ledger
updates per burst.  The emission mechanism, frequency, and
dispersion properties are all consistent with the RS framework.
FRBs serve as cosmological probes of the intergalactic medium
within RS cosmology.

%% ============================================================
\begin{thebibliography}{99}

\bibitem{Lorimer2007}
D.~R.~Lorimer \textit{et al.}, ``A Bright Millisecond Radio Burst
of Extragalactic Origin,'' Science \textbf{318}, 777 (2007).

\bibitem{Petroff2019}
E.~Petroff, J.~W.~T.~Hessels, and D.~R.~Lorimer, ``Fast Radio
Bursts,'' Astron.\ Astrophys.\ Rev.\ \textbf{27}, 4 (2019).

\bibitem{CHIME2020}
CHIME/FRB Collaboration, ``A Bright Millisecond-Duration Radio
Burst from a Galactic Magnetar,'' Nature \textbf{587}, 54 (2020).

\bibitem{RS_Axioms}
J.~Washburn, ``Recognition Science: Derivation of the Cost
Functional and Forcing Chain,'' Recognition Science Research
Institute Technical Report (2025).

\end{thebibliography}

\end{document}
